\subsection{Minimizing Aleatoric Uncertainty of Traversability Cost}

\textbf{Experimental Setup}
We quantitatively evaluate our meta-learning method for traversability cost prediction~(\textbf{Q1}). We collect driving data in various types of terrain with our vehicle~(See Fig.~\ref{fig:concept}) equipped with OS1-128 LiDAR and IMU. Also, the driving data is obtained in a simulation environment consisting of randomly patterned rough terrain and bumps. Approximately five hours of driving data are utilized to train the network. For comparison, a global model~(\textit{Baseline}) is simply trained without adaptation.

\textbf{Experimental Result}
The Table~\ref{tab:quantitative} provides the Mean Square Error~(MSE) between the ground-truth traversability and the predicted traversability cost. In every category, our method produces superior results compared to the baseline. It means the non-meta-learned baseline failed to converge due to high aleatoric uncertainty in ground-truth traversability gathered in various terrains, whereas our meta-learned model can converge well by incorporating such uncertainty during training. Our method even shows better performance without adaptation during the inference~(\textit{w.o. Adaptation}), meaning that it finds a better initial parameter. Moreover, the performance improves as the model adapts using the meta-train data derived from past interaction experiences~(\textit{w. Adaptation}). 